\documentclass[11pt]{article}
\usepackage[a4paper,margin=0.85in]{geometry}
\usepackage{amsmath,amssymb,amsthm,hyperref}
\hypersetup{hidelinks}
\newtheorem{theorem}{Theorem}
\newtheorem{lemma}{Lemma}
\newtheorem{corollary}{Corollary}
\title{Partition coefficients for generalized peakon transfer matrices}
\author{Research proof and Lean formalization}
\date{October 27, 2026}
\begin{document}
\maketitle
\begin{abstract}
We give all partition coefficients in the trace powers of a generalized peakon Lax matrix. A rank-two determinant calculation reduces its characteristic determinant to one Hamiltonian generating polynomial. The formal trace-log identity then gives a closed factorial formula, including the two coefficient families proposed by Avan, Frappat and Ragoucy. The proof applies to arbitrary coherent complex particle amplitudes and, literally, to the square-root matrices for nonnegative real momenta.
\end{abstract}
\noindent\textbf{Scope and review status.} The mathematical and formal statements below use the explicitly stated amplitude domains. Independent preliminary source and mathematical review confirms the coherent square-root convention used below. Complete final formal, manuscript and wider novelty review is pending; no priority claim or discovery credit is made here.

Let $c_i\in\mathbb C$, $q_i\in\mathbb R$ ($1\leq i\leq N$), and $\nu,\gamma,z\in\mathbb C$, with $2z+\gamma\ne0$. Define the actual hyperbolic matrices
\begin{align*}
 A_{ij}&=c_ic_j\sinh\!\left(\frac\nu2(q_i-q_j)\right),&
 T_{ij}&=c_ic_j\cosh\!\left(\frac\nu2(q_i-q_j)\right),&
 S_{ij}&=\operatorname{sgn}(q_i-q_j)A_{ij},\\
 \Lambda&=\frac{z^2-1}{2z+\gamma},&
 \xi&=\frac{z^2+\gamma z+1}{2z+\gamma},&
 L&=S+\Lambda T+\xi A.
\end{align*}
Here $\operatorname{sgn}(0)=0$. Put $\tau_m=\operatorname{tr}(L^m)$ and, exactly as in \cite[eq.~(5.17)]{AFR},
\begin{equation}\label{h}
 h^{(k)}=k\operatorname{tr}\left((T+\tfrac\gamma2 A)(S+A)^{k-1}\right)\qquad(k\geq1).
\end{equation}

\begin{theorem}\label{main}
For every $m\geq1$,
\begin{equation}\label{expansion}
 \tau_m=\sum_{j=1}^m\Lambda^j
 \sum_{\substack{k_1\geq\cdots\geq k_j\geq1\\k_1+\cdots+k_j=m}}
 \frac{m(j-1)!}{(\prod_{i=1}^j k_i)(\prod_{s\geq1}r_s!)}
 \prod_{i=1}^j h^{(k_i)},
\end{equation}
where $r_s=\#\{i:k_i=s\}$. In particular, the proposed coefficient families are
\begin{equation}\label{special}
 a_1(m)=1,\qquad a_m(1,\ldots,1)=1,\qquad
 a_{j+1}(m-j,\underbrace{1,\ldots,1}_{j})=\frac m{m-j}\quad(0\leq j<m-1).
\end{equation}
\end{theorem}
No independence or uniqueness of these monomials for fixed $N$ is asserted; the displayed universal coefficients give a valid expansion for every $N$.

\begin{corollary}\label{positive}
For arbitrary $p_i\geq0$, Theorem~\ref{main} holds for the literal factors $\sqrt{p_ip_j}$ in place of $c_ic_j$, with the same $h^{(k)}$ and all real particle positions, including ties.
\end{corollary}
Indeed choose $c_i=\sqrt{p_i}$; then $\sqrt{p_ip_j}=\sqrt{p_i}\sqrt{p_j}$. 

\section*{The determinant reduction}
First reorder the positions so that $q_1\geq\cdots\geq q_N$, and set
\[
 u_i=c_i e^{\nu q_i/2},\qquad v_i=c_i e^{-\nu q_i/2},\qquad B=S+A.
\]
Expanding $\sinh$ and $\cosh$ gives
\[
 A=\tfrac12(uv^t-vu^t),\qquad T=\tfrac12(uv^t+vu^t),\qquad
 B_{ij}=\begin{cases}u_iv_j-v_iu_j,&i<j,\\0,&i\geq j.\end{cases}
\]
These identities remain valid when positions tie, since then $A_{ij}=0$. Thus $B^N=0$, $\det(I-tB)=1$, and
\[
 R=(I-tB)^{-1}=\sum_{k=0}^{N-1}t^kB^k,\qquad B-B^t=uv^t-vu^t.
\]
All equalities are polynomial identities over a characteristic-zero field. Write
\[
 a=\frac{\Lambda+\xi-1}{2},\quad b=\frac{\Lambda-\xi+1}{2},\quad
 L=B+a uv^t+bvu^t,
\]
and put $P=v^tRu$, $Q=u^tRv$, $C=PQ-(v^tRv)(u^tRu)$. The classical rank-two determinant lemma gives
\begin{equation}\label{det2}
 \det(I-tL)=1-t(aP+bQ)+t^2abC.
\end{equation}
The resolvent identity
\[
 R-R^t=tR(B-B^t)R^t
\]
gives $P-Q=tC$ after multiplication by $v^t$ and $u$. Also
\[
 \xi^2=\Lambda^2+\gamma\Lambda+1,\qquad a+b=\Lambda,\qquad
 ab=\frac{\xi-1}{2}-\frac{\gamma\Lambda}{4}.
\]
Substitution in \eqref{det2} cancels the $\xi-1$ terms and proves
\begin{equation}\label{collapse}
 \det(I-tL)=1-\Lambda H(t),\qquad
 H(t)=t\operatorname{tr}\left((T+\tfrac\gamma2 A)R\right)
      =\sum_{k=1}^N\frac{h^{(k)}}k t^k.
\end{equation}
This is the source-specific reduction. All determinant, trace and Hamiltonian expressions are invariant under simultaneous permutation of the particle indices, so \eqref{collapse} holds without an ordering assumption. Nilpotence also proves $h^{(k)}=0$ for $k>N$.

\section*{Every trace order and every partition}
The classical formal trace-log identity, with constant determinant coefficient one, is
\[
 -\log\det(I-tL)=\sum_{m\geq1}\frac{\operatorname{tr}(L^m)}m t^m.
\]
For clarity, it can be proved by constructing the matrix geometric series $(I-tL)^{-1}$ and applying Jacobi's determinant derivative identity. Both sides then have the same derivative and zero constant term. Combining this identity with \eqref{collapse}, using $H(0)=0$, yields
\begin{equation}\label{log}
 \sum_{m\geq1}\frac{\tau_m}m t^m
   =\sum_{j\geq1}\frac{\Lambda^j}{j}H(t)^j.
\end{equation}
Only $j\leq m$ contributes to the coefficient of $t^m$. A fixed partition with multiplicities $r_s$ has $j!/\prod_s r_s!$ orderings in $H(t)^j$. Its contribution is therefore
\[
 \frac mj\,\frac{j!}{\prod_s r_s!}\prod_{i=1}^j\frac{h^{(k_i)}}{k_i},
\]
which proves \eqref{expansion}. The three evaluations in \eqref{special} follow by direct cancellation. The restriction $j<m-1$ ensures $m-j>1$; the all-ones case has its own multiplicity factor.

\section*{Source relation and formal scope}
The Lax matrices originate in \cite{RB,AFR2023}. Avan--Frappat--Ragoucy \cite[Conjecture~5.2, eqs.~(5.20)--(5.21)]{AFR} propose the partition expansion and the families \eqref{special}, leaving the other coefficients undetermined. Theorem~\ref{main} supplies all coefficients for the coherent-amplitude family; Corollary~\ref{positive} uses their displayed square-root factors literally for $p_i\geq0$. The source's notation $\sqrt{p_ip_j}$ uses coherent per-particle roots: its identity $T_{ii}=p_i$, eq.~(5.6), and its degree-$m$ homogeneity require this convention. Thus $c_i^2=p_i$ recovers its family, including complex coherent roots for negative momenta. Arbitrary $c_i$ in Theorem~\ref{main} covers this convention; the corollary additionally proves the nonnegative principal-root reading directly. The rank-two determinant lemma, Jacobi identity, formal trace-log identity and multinomial theorem are classical ingredients, not separate discoveries.

The Lean development has 121 declarations in 19 modules, using Lean 4.34.1 and mathlib revision\newline
\texttt{d13f23b723b8a846827a245b89c10fc7d3f11612}. It proves the determinant collapse, the formal trace-log and multinomial bridges, all-order support truncation, the actual complex hyperbolic matrix identities, permutation sorting, and both coefficient families. The closed source endpoint and its literal nonnegative specialization are:
\begin{center}\small\texttt{source\_conjecture\_372008}\\\texttt{nonnegative\_source\_conjecture\_372008}\end{center}
The isolated strict replay passes with only \texttt{propext}, \texttt{Classical.choice} and \texttt{Quot.sound}. Independent final formal/manuscript review is pending. Discovery credit remains zero; no global priority or optimality claim is made.

\small
\begin{thebibliography}{9}
\bibitem{AFR} J.~Avan, L.~Frappat and E.~Ragoucy, \emph{Quadratic Poisson brackets for the Camassa--Holm peakons}, J.~Phys.~A: Math.~Theor. \textbf{59} (2026), 245202. \href{https://doi.org/10.1088/1751-8121/ae7683}{doi:10.1088/1751-8121/ae7683}; arXiv:2512.11066v3.
\bibitem{AFR2023} J.~Avan, L.~Frappat and E.~Ragoucy, \emph{A new integrable structure associated to the Camassa--Holm peakons}, SciPost Phys. \textbf{15} (2023), 228. \href{https://doi.org/10.21468/SciPostPhys.15.6.228}{doi:10.21468/SciPostPhys.15.6.228}.
\bibitem{RB} O.~Ragnisco and M.~Bruschi, \emph{Peakons, R-matrix and Toda lattice}, arXiv:solv-int/9509012v1 (1995).
\end{thebibliography}
\end{document}
